% Generated draft: requires mathematical and source review.
\documentclass[11pt]{article}
\usepackage[margin=25mm]{geometry}
\usepackage{fontspec}
\setmainfont{Latin Modern Roman}
\usepackage{amsmath,amssymb,mathtools,unicode-math}
\setmathfont{Latin Modern Math}
\usepackage{enumitem,needspace,xurl,hyperref}
\setlength{\parindent}{0pt}
\setlength{\parskip}{5pt}
\setlength{\emergencystretch}{3em}
\newcommand{\N}{\mathbb{N}}
\newcommand{\Z}{\mathbb{Z}}
\newcommand{\Q}{\mathbb{Q}}
\newcommand{\R}{\mathbb{R}}
\newcommand{\C}{\mathbb{C}}
\newcommand{\F}{\mathbb{F}}
\newcommand{\A}{\mathbb{A}}
\newcommand{\PP}{\mathbb{P}}
\newcommand{\E}{\mathbb{E}}
\newcommand{\eps}{\varepsilon}
\newcommand{\dd}{\,\mathrm{d}}
\begin{document}
\section*{UnsolvedMath ID: 701016}
\textbf{Source problem code:} catalogue entry
\subsection{Definitions and mathematical statement}
Use the standard convention that a Fitting class is a nonempty class \(\mathcal X\) of finite soluble groups, closed under isomorphism, normal subgroups, and normal products. Here a group is soluble if its derived series terminates at the identity subgroup. Normal-subgroup closure means that \(N\trianglelefteq H\in\mathcal X\) implies \(N\in\mathcal X\). Normal-product closure means that, for every finite soluble group \(H\) and normal subgroups \(A,B\trianglelefteq H\), membership \(A,B\in\mathcal X\) implies \(AB\in\mathcal X\).

For a finite soluble group \(G\), define its generated Fitting class by
\[
\operatorname{Fit}(G)=\bigcap\{\mathcal X:\mathcal X\text{ is a Fitting class and }G\in\mathcal X\}.
\]
Let
\[
\mathcal L_G=\{\mathcal X:\mathcal X\text{ is a Fitting class and }\mathcal X\subseteq\operatorname{Fit}(G)\},
\]
ordered by inclusion. Classes may be identified with sets of isomorphism types. In this lattice, the meet is intersection and the join is the smallest Fitting class containing the union.

The question attributed to N.\ T.\ Vorob'\"ev is whether
\[
\text{for every finite soluble group }G,\qquad |\mathcal L_G|<\infty.
\]
The source does not explicitly define its Fitting-class conventions; the formulation above uses the nonempty, finite-soluble-group convention. Allowing the empty class would add a least element and would not change the finiteness question. The supplied registry labels the problem open; this is a supplied status, not an independently verified literature assessment.
\subsection{Short English statement}
Does every finite soluble group generate a Fitting class having only finitely many Fitting subclasses? Subclasses are ordered by inclusion, and finiteness refers to their number, not to the number of groups in the generated class.
\subsection{Sources}
\begin{itemize}
\item Local UnsolvedMath registry, permanent ID \(701016\), exact source code \texttt{catalogue entry}, title \emph{Kourovka Notebook Problem 14.31}. The supplied record attributes the problem to N.\ T.\ Vorob'\"ev, gives the proposed year as 1999, and labels it open.
\item \emph{The Kourovka Notebook (Unsolved Problems in Group Theory)}, original problem 14.31, 14th issue (1999); the supplied background also identifies the 21st-issue edition (2026), \texttt{arXiv:1401.0300v48}. The supplied collection description names E.\ I.\ Khukhro and V.\ D.\ Mazurov as editors. This bibliographic attribution is supplied source data and has not been independently checked.
\end{itemize}
\end{document}
